% ── Section: Background ──────────────────────────────────────────────────────
\section{Background}
\label{sec:background}

\noindent\textbf{ANN indexes and HNSW.}
Approximate nearest-neighbor (ANN) search finds the $k$ most similar vectors
to a query from a large database, trading a small loss in accuracy for orders
of magnitude speedup over exhaustive search.
The current state-of-the-art index is \emph{Hierarchical Navigable Small
Worlds} (HNSW)~\cite{malkov2020efficient}, a multi-layer proximity graph.
The top layers provide long-range routing; the bottom layer connects every
vector to its $M$ nearest neighbors.
A greedy beam search descends the layers, maintaining a priority queue
(``beam'') of the $\mathit{ef}$ most promising candidates.
The dominant cost is the \emph{distance computation}: for 768-dimensional
float32 vectors, computing cosine similarity requires 768 multiply-accumulate
operations ($\approx\!30$~ns on modern CPUs).
Tuning the beam width $\mathit{ef}$ provides a smooth recall--throughput
trade-off.

\smallskip
\noindent\textbf{Filtered ANN: prior approaches.}
FilteredDiskANN~\cite{gollapudi2023filtered} builds per-label entry-point
structures for efficient multi-label filtered search on disk.
ACORN~\cite{patel2024acorn} augments each node's neighbor list with additional
connections to filter-satisfying nodes, maintaining connectivity under arbitrary
predicates.
SIEVE~\cite{zhang2025sieve} (VLDB 2025) achieves high recall via $|\text{roles}|$
separate HNSW indexes, one per access role.
VBASE~\cite{zhang2023vbase} enables hybrid SQL+vector queries via
``relaxed monotonicity,'' observing that beam traversal may still make
progress through filter-failing nodes---an insight we formalize for RBAC.

\smallskip
\noindent\textbf{RBAC in healthcare.}
Role-Based Access Control~\cite{ferraiolo2001proposed} (NIST standard)
assigns users to \emph{roles}, each role carrying a set of \emph{permissions}.
In HIPAA-governed systems, roles encode department membership, clinical
function, data sensitivity categories (e.g., mental health, HIV status,
genomic data), patient consent flags, and research enrollment.
A single 64-bit bitmask can represent any conjunctive subset of these
attributes (Section~\ref{sec:dataset}), enabling the gate check
$(\mathit{node\_mask} \;\&\; \mathit{query\_mask}) = \mathit{query\_mask}$
as a single bitwise AND instruction.
